\begin{table}[ht]
\centering
\caption{Recovery time and pod restarts during fault injection (median, n=30)}
\label{tab:overhead}
\small
\begin{tabular}{clrrrr}
\hline
\textbf{ID} & \textbf{Scenario} & \multicolumn{2}{c}{\textbf{Recovery time (s)}} & \multicolumn{2}{c}{\textbf{Pod restarts}} \\
& & CM & LT & CM & LT \\
\hline
P1 & Pod Kill & 0.0 & 0.0 & 0 & 0 \\
P2 & Container Kill & 0.0 & 0.0 & 1 & 3 \\
P3 & Pod Failure & 0.0 & 0.0 & 3 & 0 \\
N1 & Latency 50ms & 0.0 & 0.0 & 0 & 0 \\
N2 & Latency 100ms & 0.0 & 0.0 & 0 & 0 \\
N3 & Latency 300ms & 0.0 & 0.0 & 0 & 0 \\
N4 & Packet Loss 5\% & 0.0 & 0.0 & 0 & 0 \\
N5 & Network Partition & 0.0 & 0.0 & 0 & 0 \\
R1 & CPU Stress 80\% & 0.0 & 0.0 & 0 & 0 \\
R2 & Memory Pressure 80\% & 0.0 & 0.0 & 0 & 0 \\
A1 & HTTP Abort 503 & 0.0 & 0.0 & 0 & 0 \\
A2 & gRPC Unavailable & 0.0 & 0.0 & 0 & 0 \\
\hline
\end{tabular}
\vspace{2mm}
\raggedright\footnotesize Recovery time = elapsed seconds from fault end until mean CPU across faulted pods first returns to within 20\% of that run's own baseline-phase CPU mean, right-censored at 60s (this study's recovery-phase length) if never reached; an operational definition introduced for this analysis, not literally pre-specified beyond naming ``recovery time'' as a metric (see PREREGISTRATION.md and analysis/analyze.py docstring).
\end{table}